\documentclass[11pt]{article}

\usepackage[margin=1in]{geometry}
\usepackage{amsmath,amssymb,amsthm}
\usepackage{booktabs}
\usepackage{enumitem}
\usepackage{hyperref}
\usepackage{xcolor}

\title{Lecture~8 Study Notes:\\Transformers, Part~1}
\author{Deep Learning Notes\\
\small Ground-up intuitions and worked examples\\[0.35em]
\small Formation borrows heavily from the
\href{https://r-three.github.io/deep-learning-lecture-notes/lectures/08-transformers-part-1.html}{original course website (CSC2516)}}
\date{}

\begin{document}
\maketitle

\begin{abstract}
These notes are a first slice of \textbf{Lecture~8: Transformers, Part~1}.
Still to come: removing biases, gated feed-forward networks, and the full
BERT write-up. What is here so far is \textbf{layer normalization} set next
to batch normalization, the \textbf{position-wise feed-forward network},
\textbf{pre-} versus \textbf{post-layer norm}, \textbf{efficient and local
attention}, \textbf{grouped-query attention}, and the \textbf{encoder-only
transformer}---with application intuition for classification, labeling,
span extraction, and fine-tuning versus frozen embeddings---with the baby
examples kept for reference. Content formation borrows heavily
from the original
\href{https://r-three.github.io/deep-learning-lecture-notes/}{course website
(CSC2516)}.
\end{abstract}

\tableofcontents
\vspace{1em}

%----------------------------------------------------------------------
\section{Layer Normalization}
%----------------------------------------------------------------------

Self-attention, position encodings, and multi-head attention are already in
the Lecture~7 notes. The new piece in the transformer is layer
normalization \cite{ba2016layer}. It plays the same role batch normalization
played in a convolutional network: keep the scale of activations from
drifting as they pass through the model. The calculation uses a different
group of numbers.

\noindent\textbf{One-line picture.}
Batch norm curves one quiz using every student. Layer norm curves one
student using every quiz. In a transformer, the student is one word and the
quiz scores are the entries of that word's vector.

\paragraph{Definition.}
Given a length-\(T\) sequence \(\mathbf{h}_1,\ldots,\mathbf{h}_T\in\mathbb{R}^d\),
layer normalization computes, at each position \(t\), the mean and standard
deviation over that position's \(d\) features,
\[
  \mu_t
  =
  \frac{1}{d}\sum_{i=1}^{d} h_{t,i},
  \qquad
  \sigma_t
  =
  \sqrt{\frac{1}{d}\sum_{i=1}^{d}(h_{t,i}-\mu_t)^2},
\]
and outputs
\[
  \operatorname{layernorm}(\mathbf{h})_t
  =
  \boldsymbol{\gamma}
  \odot
  \frac{\mathbf{h}_t-\mu_t}{\sigma_t}
  +
  \boldsymbol{\beta}.
\]
The vectors \(\boldsymbol{\gamma}\) and \(\boldsymbol{\beta}\) are learned.
They let the model stretch and shift the normalized vector, so the layer is
allowed to leave the ``mean \(0\), spread \(1\)'' result behind if that is
too strict. In every baby example below they are \(1\) and \(0\), so the
normalized numbers are the final output.

Batch normalization, for comparison, standardizes each feature using the
mean and variance of the current batch and then applies the same kind of
learned scale and shift \cite{ioffe2015batch}:
\[
  \operatorname{BN}(\mathbf{x})
  =
  \boldsymbol{\gamma}
  \odot
  \frac{\mathbf{x}-\boldsymbol{\mu}_B}{\boldsymbol{\sigma}_B+\epsilon}
  +
  \boldsymbol{\beta}.
\]
Here \(\boldsymbol{\mu}_B\) and \(\boldsymbol{\sigma}_B\) run over the
examples in the batch, one number per feature, and \(\epsilon\) is a small
constant that keeps the division safe when the spread is tiny. The baby
arithmetic below uses spreads of exactly \(1\), so \(\epsilon\) does not
change the displayed digits.

%----------------------------------------------------------------------
\section{What \(h_{t,i}\) Is}
\label{sec:hti}
%----------------------------------------------------------------------

A sequence is a list of word vectors. Stack the words as the rows of a
table. The row number is the position \(t\). The column number is the
feature \(i\). The symbol \(h_{t,i}\) is the single number sitting in row
\(t\), column \(i\).

Take the two-word sequence ``coffee cold.'' Each word is a vector of length
\(d=2\), so the sequence length is \(T=2\):

\begin{center}
\begin{tabular}{@{}lcc@{}}
\toprule
 & feature \(i=1\) & feature \(i=2\) \\
\midrule
\(t=1\), coffee & \(h_{1,1}=0\) & \(h_{1,2}=2\) \\
\(t=2\), cold & \(h_{2,1}=2\) & \(h_{2,2}=4\) \\
\bottomrule
\end{tabular}
\end{center}

The whole vector for a word is the row. Coffee is
\(\mathbf{h}_1 = [h_{1,1},\, h_{1,2}] = [0,\; 2]\).
Cold is \(\mathbf{h}_2 = [2,\; 4]\).
One entry is only one cell: \(h_{1,1}=0\) is coffee's first feature, and
\(h_{2,2}=4\) is cold's second feature.

The sum in the layer-norm formula adds across one row, with \(t\) held
still and \(i\) running:
\[
  \sum_{i=1}^{d} h_{1,i}
  =
  h_{1,1}+h_{1,2}
  =
  0+2.
\]
That is coffee only. Cold's row is a separate sum,
\(h_{2,1}+h_{2,2}=2+4\).

%----------------------------------------------------------------------
\section{Mean over Axis 1 instead of Axis 0}
\label{sec:axes}
%----------------------------------------------------------------------

The same table has two directions.

\begin{itemize}[leftmargin=*]
  \item \textbf{Axis 0} points down the rows: from coffee to cold. In a
        batch of separate examples this is the batch direction. On this
        sentence the rows are the two positions, so axis 0 points down the
        sequence.
  \item \textbf{Axis 1} points across a row: from feature \(1\) to feature
        \(2\) inside one word.
\end{itemize}

A mean has to choose one direction to collapse.

\textbf{Mean over axis 0} walks down a column. Feature \(1\) is the column
\([0,\; 2]\), and its mean is \(1\). Feature \(2\) is the column
\([2,\; 4]\), and its mean is \(3\). Coffee and cold are averaged together.
That is the batch-norm direction: one feature, pooled over every row.

\textbf{Mean over axis 1} walks across a row. Coffee's row \([0,\; 2]\) has
mean \(1\). Cold's row \([2,\; 4]\) has mean \(3\). The other word is not
in the average. That is layer norm, and it is the sum written above divided
by \(d=2\).

Both spreads on this table equal \(1\), so dividing by the spread gives:

\begin{center}
\begin{tabular}{@{}lcc@{}}
\toprule
 & after a mean on axis 0 & after a mean on axis 1 \\
 & (batch-norm direction) & (layer norm) \\
\midrule
coffee & \([-1,\; -1]\) & \([-1,\; 1]\) \\
cold & \([1,\; 1]\) & \([-1,\; 1]\) \\
\bottomrule
\end{tabular}
\end{center}

The line ``we compute a mean over axis 1 instead of axis 0'' is this swap
and nothing more. Batch norm's mean runs down the columns. Layer norm uses
the same table and runs the mean across each row instead. In the formula,
that is a mean over \(i\) at a fixed \(t\), which is a mean over
\(h_{t,1}, h_{t,2}, \ldots, h_{t,d}\).

%----------------------------------------------------------------------
\section{One Word, Four Features}
%----------------------------------------------------------------------

Take the word ``cold'' inside the sentence ``the coffee was cold,'' and
suppose the vector that reaches layer norm is
\[
  \mathbf{h} = [2,\; 4,\; 4,\; 6].
\]
Those four numbers are four learned features of that one word.

\[
  \mu = \frac{2+4+4+6}{4} = 4,
  \qquad
  \sigma
  =
  \sqrt{\frac{(-2)^2+0^2+0^2+2^2}{4}}
  =
  \sqrt{2}.
\]
\[
  \frac{\mathbf{h}-\mu}{\sigma}
  =
  [-\sqrt{2},\; 0,\; 0,\; \sqrt{2}]
  \approx
  [-1.41,\; 0,\; 0,\; 1.41].
\]
``Coffee'' is normalized from its own four numbers. The other words in the
sentence do not enter, and neither do the other sentences in the batch.

%----------------------------------------------------------------------
\section{Baby Example: Two Quizzes}
%----------------------------------------------------------------------

The same rule is easier to see with two students and two scores. Spelling
is one feature; math is the other.

\begin{center}
\begin{tabular}{@{}lcc@{}}
\toprule
Student & Spelling & Math \\
\midrule
A & 0 & 2 \\
B & 2 & 4 \\
\bottomrule
\end{tabular}
\end{center}

\subsection{Batch norm}

Look down one column at a time, across the class.

Spelling scores are \(0\) and \(2\). Their average is \(1\), and each score
sits exactly \(1\) away from that average:
\[
  \text{A: }\frac{0-1}{1} = -1,
  \qquad
  \text{B: }\frac{2-1}{1} = 1.
\]
Math scores are \(2\) and \(4\). Average \(3\), same spread:
\[
  \text{A: }\frac{2-3}{1} = -1,
  \qquad
  \text{B: }\frac{4-3}{1} = 1.
\]

\begin{center}
\begin{tabular}{@{}lcc@{}}
\toprule
Student & Spelling & Math \\
\midrule
A & \(-1\) & \(-1\) \\
B & \(1\) & \(1\) \\
\bottomrule
\end{tabular}
\end{center}

A's spelling grade moved because of B's spelling grade. A different
classmate in the batch would change A's curved spelling score, even if A
wrote the same answers.

\subsection{Layer norm}

Look across one row at a time, using only that student's own scores.

Student A has \(0\) and \(2\). Average \(1\), spread \(1\):
\[
  \frac{0-1}{1} = -1,
  \qquad
  \frac{2-1}{1} = 1.
\]
Student B has \(2\) and \(4\). Average \(3\), spread \(1\):
\[
  \frac{2-3}{1} = -1,
  \qquad
  \frac{4-3}{1} = 1.
\]

\begin{center}
\begin{tabular}{@{}lcc@{}}
\toprule
Student & Spelling & Math \\
\midrule
A & \(-1\) & \(1\) \\
B & \(-1\) & \(1\) \\
\bottomrule
\end{tabular}
\end{center}

Each student is compared only with themselves. B could be absent and A's
result would stay \(-1\) and \(1\).

%----------------------------------------------------------------------
\section{The Same Numbers as Word Vectors}
%----------------------------------------------------------------------

Carry the quiz table into a transformer. Each word is one student. The two
numbers inside its vector are its two quiz scores. A batch of two one-word
sentences arrives at the normalization:

\begin{center}
\begin{tabular}{@{}llc@{}}
\toprule
Sentence & Word & Vector \\
\midrule
1 & coffee & \([0,\; 2]\) \\
2 & tea & \([2,\; 4]\) \\
\bottomrule
\end{tabular}
\end{center}

The two entries are two learned features of the word, for instance ``how
drink-like'' and ``how hot or cold.''

\subsection{Layer norm, the one a transformer uses}

Each word is rescaled using only its own two numbers.

\textbf{coffee} \(= [0,\; 2]\). Average \((0+2)/2 = 1\), spread \(1\):
\[
  \frac{0-1}{1} = -1,
  \qquad
  \frac{2-1}{1} = 1,
  \qquad
  \text{coffee becomes } [-1,\; 1].
\]

\textbf{tea} \(= [2,\; 4]\). Average \((2+4)/2 = 3\), spread \(1\):
\[
  \frac{2-3}{1} = -1,
  \qquad
  \frac{4-3}{1} = 1,
  \qquad
  \text{tea becomes } [-1,\; 1].
\]

\begin{center}
\begin{tabular}{@{}lcc@{}}
\toprule
Word & Before & After layer norm \\
\midrule
coffee & \([0,\; 2]\) & \([-1,\; 1]\) \\
tea & \([2,\; 4]\) & \([-1,\; 1]\) \\
\bottomrule
\end{tabular}
\end{center}

``Tea'' never enters the arithmetic for ``coffee.'' In a longer sentence
such as ``the coffee was cold,'' the same thing happens separately for
every word. Inside the transformer block, this is the table that runs:
one word vector in, one rescaled word vector out.

\subsection{Batch norm on these same vectors}

Batch norm takes one feature column down the batch.

First column, the first number of each word: \(0\) and \(2\). Average \(1\),
spread \(1\).
\[
  \text{coffee: }\frac{0-1}{1} = -1,
  \qquad
  \text{tea: }\frac{2-1}{1} = 1.
\]
Second column: \(2\) and \(4\). Average \(3\), spread \(1\).
\[
  \text{coffee: }\frac{2-3}{1} = -1,
  \qquad
  \text{tea: }\frac{4-3}{1} = 1.
\]

\begin{center}
\begin{tabular}{@{}lcc@{}}
\toprule
Word & Before & After batch norm \\
\midrule
coffee & \([0,\; 2]\) & \([-1,\; -1]\) \\
tea & \([2,\; 4]\) & \([1,\; 1]\) \\
\bottomrule
\end{tabular}
\end{center}

Here ``coffee'' changes because ``tea'' is sitting in the same batch.

%----------------------------------------------------------------------
\section{Where Batch Normalization Is Used}
%----------------------------------------------------------------------

Batch normalization is the tool from Lecture~5, inside image models, and in
particular inside a ResNet block \cite{he2016deep}. The order in one block is
\begin{enumerate}[leftmargin=*, nosep]
  \item \(3\times 3\) convolution, batch norm, ReLU;
  \item \(3\times 3\) convolution, batch norm;
  \item add the residual, then ReLU.
\end{enumerate}
The network also opens with a large \(7\times 7\) convolution, then batch
norm, then max pooling, and then a stack of those blocks.

On an image, one activation is a volume: a batch of pictures, each with
height, width, and channels. Batch norm keeps a separate mean and spread
for each channel, and it computes them from every picture in the batch and
every pixel in that channel. A ``red-edge'' channel is curved using all the
red-edge responses in the batch, the same way the spelling quiz was curved
using every student's spelling score.

At test time there may be no batch. The model then uses a running average
of the means and spreads it collected during training.

The quiz table already shows the uncomfortable part: one example's output
depends on the other examples in the batch. Lecture~5 records that
dependence, and later models put layer norm in the corresponding slot. In
the transformer, that slot is the per-word calculation above.

%----------------------------------------------------------------------
\section{The Position-Wise Feed-Forward Network}
\label{sec:ffn}
%----------------------------------------------------------------------

This follows the course lecture's account of the transformer's ingredients
(\href{https://r-three.github.io/deep-learning-lecture-notes/lectures/08-transformers-part-1.html}{Lecture~8}).
Inside a block, self-attention is the part that lets words look at each
other. The feed-forward piece does a different job: it rewrites each word
on its own.

The original transformer uses one hidden layer and a ReLU
\cite{vaswani2017attention},
\[
  \mathbf{h}_t
  =
  \mathbf{W}_2\,
  \phi(\mathbf{W}_1 \mathbf{x}_t + \mathbf{b}_1)
  +
  \mathbf{b}_2.
\]
\(\mathbf{W}_1\) maps the \(d\)-dimensional word up to a wider hidden
vector, \(\phi\) is applied entry by entry, and \(\mathbf{W}_2\) maps back
to \(d\) dimensions. The same \(\mathbf{W}_1\), \(\mathbf{W}_2\),
\(\mathbf{b}_1\), and \(\mathbf{b}_2\) are used at every position, the same
sharing pattern as \(\boldsymbol{\gamma}\) and \(\boldsymbol{\beta}\): \(T\)
applications, one shared set of weights.

\subsection{Why it is pointwise}

``Pointwise'' means the function looks at one position and ignores the
others. On the sentence ``coffee cold,'' the network reads
\([0,\; 2]\) by itself and \([2,\; 4]\) by itself. Take
\[
  \mathbf{W}_1
  =
  \begin{pmatrix} 1 & -1 \\ 0 & 1 \end{pmatrix},
  \qquad
  \mathbf{b}_1=\mathbf{0},
  \qquad
  \phi=\operatorname{ReLU},
  \qquad
  \mathbf{W}_2=\mathbf{I},
  \qquad
  \mathbf{b}_2=\mathbf{0}.
\]
Coffee \([0,\; 2]\) is sent to \([-2,\; 2]\), and ReLU zeroes the negative
entry, leaving \([0,\; 2]\). Cold \([2,\; 4]\) is sent to \([-2,\; 4]\),
and ReLU leaves \([0,\; 4]\). Cold's output was computed without ever
reading coffee's numbers. Swap the two words and the two outputs swap with
them. A third word would be passed through these same two matrices; the
parameter count does not grow with \(T\).

That is the useful split inside the block. Attention is where ``cold'' can
pick up information from ``coffee.'' The feed-forward network is where each
word, already mixed, is rewritten by a small MLP. Layer norm is the same
kind of per-word step: shuffle the input words and the outputs shuffle in
the same order.

\subsection{What is different from a traditional MLP}

The recipe inside one position is an ordinary MLP: linear layer,
nonlinearity, linear layer. The difference is what counts as the input.

A traditional MLP for the whole sentence flattens the rows into one long
vector. ``Coffee cold'' becomes \([0,\; 2,\; 2,\; 4]\), and one weight
matrix connects every entry of that vector to every output. A single weight
can carry coffee into cold's output. With the weight \(5\) on coffee's
second feature, that feature contributes \(5\cdot 2 = 10\) to whichever
output the weight feeds, including an output that is supposed to describe
cold. The matrix also has shape tied to \(T\): a length-\(2\) sentence and
a length-\(3\) sentence are different input sizes, so they do not share one
matrix.

The position-wise network has no such cross-word weight. Each application
sees a vector of length \(d\) only, and the matrices stay \(d\) by \(d\)
(or \(d\) by the hidden width) no matter how long the sentence is.

%----------------------------------------------------------------------
\section{Pre-Layer Norm and Post-Layer Norm}
\label{sec:pre-post}
%----------------------------------------------------------------------

This is the first of the block modifications in the same course lecture.
The original transformer puts layer norm \emph{after} the residual add.
That arrangement is now called \textbf{post-layer norm}. The alternative,
\textbf{pre-layer norm}, puts layer norm \emph{before} the attention or the
feed-forward network, and the residual add then wraps around both.

Post-layer norm, the 2017 block \cite{vaswani2017attention}:
\begin{align*}
  \mathbf{h}
  &=
  \operatorname{LN}\bigl(\mathbf{x} + \operatorname{MHSA}(\mathbf{x})\bigr),
  \\
  \mathbf{y}
  &=
  \operatorname{LN}\bigl(\mathbf{h} + \operatorname{FFN}(\mathbf{h})\bigr).
\end{align*}
The skip connection is added first, and layer norm sees the sum. The copy
of \(\mathbf{x}\) that skipped the attention has to pass through the layer
norm before it leaves the block.

Pre-layer norm:
\begin{align*}
  \mathbf{h}
  &=
  \mathbf{x} + \operatorname{MHSA}\bigl(\operatorname{LN}(\mathbf{x})\bigr),
  \\
  \mathbf{y}
  &=
  \mathbf{h} + \operatorname{FFN}\bigl(\operatorname{LN}(\mathbf{h})\bigr).
\end{align*}
Layer norm is applied to the input of the attention, or of the feed-forward
network. The skip adds the original vector after that update has been
computed. Stacked from the embeddings to the output, those skips form a
path with nothing on it: no attention, no feed-forward network, and no
layer norm. The course lecture calls this the uninterrupted residual
stream.

The add is legal with no extra projection because the block is built so
that every piece preserves the width \(d_{\text{model}}\): multi-head
attention, layer norm, and the feed-forward network all map
\(\mathbb{R}^{d}\) to \(\mathbb{R}^{d}\).

\subsection{What the placement changes}

Along the pre-layer-norm skip, the gradient from the output back to the
input is the identity on that path. Deep stacks can still send a learning
signal to early blocks when the attention and the feed-forward branch are
difficult to differentiate. The course lecture is careful about the
strength of that claim: which placement trains better depends on
initialization, the optimizer, and the other hyperparameters, so it is not
a settled comparison. Pre-layer norm is the more common choice in current
models. A transformer built with the original post-layer-norm block still
works.

The other reading is about what the network is doing to the sentence. With
pre-layer norm, the running vectors start as the word embeddings. Each
block computes an update from a normalized copy and adds that update back.
The representation is revised in place, rather than being replaced by a
freshly normalized sum at every block.

%----------------------------------------------------------------------
\section{Efficient and Sub-Quadratic Attention}
\label{sec:subquadratic}
%----------------------------------------------------------------------

This follows the course lecture's section on efficient self-attention.
Full self-attention scores every word against every word. A sentence of
length \(T\) therefore builds a \(T\times T\) score matrix: \(T^2\) dot
products. For ``the coffee was cold,'' \(T=4\), that is \(16\) dots. For a
\(1000\)-word document it is \(1000^2 = 10^6\) dots in every head of every
block.

The lecture gives two reasons this quadratic bill is often left alone, and
then the simple change that actually reduces it: local self-attention.

\subsection{The feed-forward network is often the larger bill}

Count multiplies in one block, with width \(d\) and feed-forward width
\(d_{\mathrm{ff}}\). The feed-forward network applies two matrices at each
of the \(T\) positions, about \(2\,T\,d\,d_{\mathrm{ff}}\) multiplies.
Score-and-average attention, with the head widths adding up to \(d\), is
about \(2\,T^2 d\) multiplies.

Take the lecture's regime, \(d=256\) and \(d_{\mathrm{ff}}=4d=1024\):
\begin{align*}
  \text{feed-forward}
  &\approx 2\cdot T\cdot 256\cdot 1024 = 524288\,T,\\
  \text{attention}
  &\approx 2\cdot T^2\cdot 256 = 512\,T^2.
\end{align*}
At \(T=256\), the feed-forward count is four times the attention count.
The two meet near \(T\approx 1024\). Past that, attention is the larger
bill. The lecture's statement that the quadratic term is not the main cost
for modest lengths, roughly \(T<1000\), is this comparison. These counts
are the big matrix products only; they are not a full account of hardware
time.

\subsection{Flash Attention stores less, and still scores every pair}

The other quadratic cost is memory. Writing down every score, for \(L\)
blocks and \(H\) heads, is \(L\cdot H\cdot T\cdot T\) numbers. On a GPU the
expensive part is moving those numbers between the large memory and the
small on-chip memory. \textbf{Flash Attention} still computes every
query--key dot product, block by block, and keeps a running softmax so the
full \(T\times T\) matrix is never stored \cite{dao2022flashattention}. The
arithmetic stays quadratic. The stored working set does not.

One query, four keys, scores \([0,\;\ln 2,\;0,\;\ln 2]\). The full softmax
uses \(\mathrm{e}^{0}=1\) and \(\mathrm{e}^{\ln 2}=2\), so the weights are
\[
  \frac{[1,\;2,\;1,\;2]}{6}
  =
  \Bigl[\tfrac{1}{6},\;\tfrac{2}{6},\;\tfrac{1}{6},\;\tfrac{2}{6}\Bigr].
\]
The running version sees the scores in two blocks of two and only remembers
a running maximum \(m\) and a running sum \(\ell\) of
\(\mathrm{e}^{s-m}\).

Block \([0,\;\ln 2]\): \(m=\ln 2\), and
\(\ell=\mathrm{e}^{0-\ln 2}+\mathrm{e}^{0}=\tfrac{1}{2}+1=\tfrac{3}{2}\).
The second block is the same, so the new maximum is still \(\ln 2\) and the
sums add, \(\ell=\tfrac{3}{2}+\tfrac{3}{2}=3\). A weight is
\(\mathrm{e}^{s-m}/\ell\): the score \(0\) gives
\((\tfrac{1}{2})/3=\tfrac{1}{6}\), and the score \(\ln 2\) gives
\(1/3=\tfrac{2}{6}\). Same weights as the full softmax. The four scores
never had to sit in memory together. If the second block had contained a
larger score, the old \(\ell\) would be multiplied by
\(\mathrm{e}^{m_{\mathrm{old}}-m_{\mathrm{new}}}\) before the new block was
added, which is what keeps the result exact.

\subsection{Local self-attention}
\label{sec:local}

When the document really is long, the lecture's simplest cut is
\textbf{local self-attention}: position \(t\) scores only a window of width
\(K\) around itself. The score matrix is a band around the diagonal. The
dot-product count drops from \(T^2\) to about \(K T\). For
\(T=1000\) and \(K=64\) that is \(64000\) dots in place of \(10^6\).

\paragraph{Baby sentence.}
Four words, values \(\mathbf{v}=[1,\;3,\;5,\;7]\) for ``the coffee was
cold,'' and a window of radius \(1\): each word may look at itself and its
immediate neighbors. Inside this first pass every allowed score is \(0\),
so the softmax is uniform on the window. The band is
\begin{center}
\begin{tabular}{@{}lcccc@{}}
\toprule
query \(\backslash\) key & the & coffee & was & cold \\
\midrule
the    & \(1\) & \(1\) & \(0\) & \(0\) \\
coffee & \(1\) & \(1\) & \(1\) & \(0\) \\
was    & \(0\) & \(1\) & \(1\) & \(1\) \\
cold   & \(0\) & \(0\) & \(1\) & \(1\) \\
\bottomrule
\end{tabular}
\end{center}
Ten scores are computed, rather than sixteen. The outputs are averages of
the values inside each row:
\begin{align*}
  \text{the}    &= \tfrac{1+3}{2} = 2,\\
  \text{coffee} &= \tfrac{1+3+5}{3} = 3,\\
  \text{was}    &= \tfrac{3+5+7}{3} = 5,\\
  \text{cold}   &= \tfrac{5+7}{2} = 6.
\end{align*}
``Cold'' reads ``was'' and ``cold'' only. The value \(1\) on ``the'' is
absent from \(6\).

\paragraph{The weights inside the window still come from the words.}
Replace ``the'''s two scores by \([\ln 2,\;0]\). The weights become
\(\tfrac{2}{3}\) and \(\tfrac{1}{3}\), and the output becomes
\[
  \tfrac{2}{3}\cdot 1 + \tfrac{1}{3}\cdot 3 = \tfrac{5}{3},
\]
instead of the flat average \(2\). A convolution would apply one fixed
kernel at every window, with no query and no key. Local attention is a
softmax over dot products, restricted to the band. The band is what it
shares with a convolution; the content-dependent weights are the
difference.

\paragraph{Path length.}
Global attention gives every word a direct look at every other word, so
``the'' can affect ``cold'' in one layer. In the band above, ``the''
affects ``coffee'' in one layer, ``was'' only on the next layer (when
``was'' reads ``coffee'''s new vector), and ``cold'' on the layer after
that. The number of layers grows with the distance, as it does in a
convolutional stack. A common repair, stated in the lecture, is to keep
most blocks local and insert an occasional block with the full \(T\times T\)
matrix, so a long jump is available without paying for it at every layer.

%----------------------------------------------------------------------
\section{Grouped-Query Attention}
\label{sec:gqa}
%----------------------------------------------------------------------

Multi-head attention gives every head its own queries, keys, and values.
The course lecture then describes two reductions. \textbf{Multi-query
attention} keeps a separate query per head and shares one key projection
and one value projection across all heads \cite{shazeer2019fast}.
\textbf{Grouped-query attention} sits between that and ordinary multi-head
attention: \(G\) groups of keys and values, with \(G\) dividing the number
of heads \(H\), and still one query per head \cite{ainslie2023gqa}.
\(G=1\) is multi-query attention. \(G=H\) is ordinary multi-head attention.

The saving the lecture emphasizes is memory traffic, especially when past
keys and values are kept for word-by-word generation (Lecture~9). Fewer
distinct key and value vectors means fewer vectors to store and to move.
Sharing every head's keys, the \(G=1\) end, is a strong constraint and can
hurt quality slightly. A middle \(G\) keeps most of the saving and less of
that loss.

\subsection{One starting point for all three}

Two words, ``coffee'' at position \(1\) and ``cold'' at position \(2\).
Four heads, and every query at position \(1\) is \(q=1\). Keys and values
are scalars, so \(\sqrt{d_k}=1\) and a score is the product \(q\,k\).
Multi-head attention gives each head a private memory. That memory is the
starting point. Grouped-query and multi-query attention keep these same
queries and only tie some of the memories together.

\begin{center}
\small
\begin{tabular}{@{}cccc@{}}
\toprule
Head & keys (coffee, cold) & values (coffee, cold) & name \\
\midrule
1 & \([\ln 2,\; 0]\) & \([3,\; 6]\) & \(M_1\) \\
2 & \([0,\; \ln 2]\) & \([0,\; 9]\) & \(M_2\) \\
3 & \([\ln 2,\; 0]\) & \([6,\; 3]\) & \(M_3\) \\
4 & \([0,\; \ln 2]\) & \([9,\; 0]\) & \(M_4\) \\
\bottomrule
\end{tabular}
\end{center}

Head~1 against \(M_1\): scores \([1\cdot\ln 2,\; 1\cdot 0]=[\ln 2,\; 0]\),
weights \(\bigl(\tfrac{2}{3},\,\tfrac{1}{3}\bigr)\), output
\[
  \tfrac{2}{3}\cdot 3 + \tfrac{1}{3}\cdot 6 = 4.
\]
Head~2 against its own \(M_2\): scores \([0,\;\ln 2]\), weights
\(\bigl(\tfrac{1}{3},\,\tfrac{2}{3}\bigr)\), output
\[
  \tfrac{1}{3}\cdot 0 + \tfrac{2}{3}\cdot 9 = 6.
\]
Head~3 against \(M_3\) has the same keys as head~1, so the same weights
\(\bigl(\tfrac{2}{3},\,\tfrac{1}{3}\bigr)\), and different values:
\(\tfrac{2}{3}\cdot 6 + \tfrac{1}{3}\cdot 3 = 5\).
Head~4 against \(M_4\) has the same keys as head~2, so the same weights
\(\bigl(\tfrac{1}{3},\,\tfrac{2}{3}\bigr)\), and output
\(\tfrac{1}{3}\cdot 9 + \tfrac{2}{3}\cdot 0 = 3\).
Four private memories, four outputs: \(4,\;6,\;5,\;3\).

\subsection{What each method is allowed to read}

The queries stay \(q=1\) on every head. The methods differ only in which
memory a head may use.
\begin{itemize}[leftmargin=*]
  \item \textbf{Multi-head} (\(G=H=4\)): head \(h\) reads \(M_h\).
  \item \textbf{Grouped-query} (\(G=2\)): heads \(1\) and \(2\) both read
        \(M_1\); heads \(3\) and \(4\) both read \(M_3\). The memories
        \(M_2\) and \(M_4\) are dropped.
  \item \textbf{Multi-query} (\(G=1\)): all four heads read \(M_1\).
        The memories \(M_2\), \(M_3\), and \(M_4\) are dropped.
\end{itemize}
Head~2 under grouped-query attention therefore recomputes head~1's
arithmetic and outputs \(4\), where multi-head attention had given it
\(6\). Head~4 recomputes head~3's arithmetic and outputs \(5\), where
multi-head attention had given it \(3\). Multi-query attention sends heads
\(3\) and \(4\) to \(M_1\) as well, so they output \(4\) too.

\begin{center}
\small
\begin{tabular}{@{}clll@{}}
\toprule
Head & Multi-head & Grouped-query & Multi-query \\
\midrule
1 & \(M_1\to 4\) & \(M_1\to 4\) & \(M_1\to 4\) \\
2 & \(M_2\to 6\) & \(M_1\to 4\) & \(M_1\to 4\) \\
3 & \(M_3\to 5\) & \(M_3\to 5\) & \(M_1\to 4\) \\
4 & \(M_4\to 3\) & \(M_3\to 5\) & \(M_1\to 4\) \\
\bottomrule
\end{tabular}
\end{center}

Reading across a row is the contrast. Head~1 is untouched, because every
method keeps \(M_1\). Head~2 changes as soon as a method makes it share
with head~1. Heads \(3\) and \(4\) still have their own group under
grouped-query attention, and both collapse onto \(M_1\) only under
multi-query attention. The lecture's remark that sharing every key and
value is a strong assumption is this last column: one memory, and every
head emits \(4\).

For \(T=2\), the number of stored key vectors, and likewise value vectors,
is \(H\) or \(G\) per position:
\begin{center}
\begin{tabular}{@{}lcc@{}}
\toprule
 & key vectors & value vectors \\
\midrule
multi-head, \(G=4\) & \(4\cdot 2=8\) & \(8\) \\
grouped-query, \(G=2\) & \(2\cdot 2=4\) & \(4\) \\
multi-query, \(G=1\) & \(1\cdot 2=2\) & \(2\) \\
\bottomrule
\end{tabular}
\end{center}
The query side does not shrink. Four heads still mean four queries at each
position, in all three columns.

%----------------------------------------------------------------------
\section{Encoder-Only Transformers}
\label{sec:encoder-only}
%----------------------------------------------------------------------

This follows the course lecture's encoder-only section. The original
transformer is an encoder--decoder. The simplest use of the block keeps
only the encoder half: the same layer, repeated \(L\) times. A length-\(T\)
sequence goes in and a length-\(T\) sequence comes out, one vector per input
position. It cannot emit a new sentence of a different length. The lecture
uses that one-to-one shape for three supervised tasks.

\subsection{Three supervised uses}

\textbf{Sequence classification.} Append a \textbf{CLS} position, as in
``the coffee was cold \texttt{[CLS]}.'' Run the stack. The vector at
\texttt{[CLS]} is a summary of the whole sentence, the role the final
hidden state played for an RNN classifier. A linear classifier reads only
that vector. Suppose it produces two scores, complaint \(=2\) and other
\(=0\). Softmax turns them into probabilities proportional to
\(\mathrm{e}^{2}\) and \(\mathrm{e}^{0}=1\):
\[
  P(\text{complaint})
  =
  \frac{\mathrm{e}^{2}}{\mathrm{e}^{2}+1},
  \qquad
  P(\text{other})
  =
  \frac{1}{\mathrm{e}^{2}+1}.
\]
The training target is one of those two classes. For a review that a person
marked as a complaint, the target is complaint. That mark is a human label,
so this training step is supervised.

\textbf{Sequence labeling.} One label per position. On the same sentence a
person, or a tagged corpus, can supply
\[
  \text{the/DET},\ 
  \text{coffee/NOUN},\ 
  \text{was/VERB},\ 
  \text{cold/ADJ}.
\]
Each output vector is classified on its own. The loss is a sum of
per-position cross-entropies against those tags. Supervised again: the tags
are the annotation.

\textbf{Span extraction.} The input is a question followed by a context
paragraph, and the answer is a stretch of that paragraph. Each output
position is classified as the start of the answer, the end, or neither. The
lecture's picture is ``when was so-and-so born?'' in front of a sentence
that contains ``1953'': the target marks the start and the end of
``1953.'' Nothing new is written; a region of the input is named. The start
and end marks are human labels.

\subsection{Application intuition: when an encoder is enough}

The encoder-only stack is a bidirectional reader. Every output position
already saw every input position, so it is a good fit when the job is to
\emph{understand} a fixed text and attach a decision to it---not when the
job is to \emph{write} a longer answer that was never in the input.

\paragraph{One label for the whole document.}
News topic classification is the running picture. A headline and a short
blurb---``Ford to Offer Sirius on More Vehicles'' plus the body---should
map to Business, not Sports. That is exactly sequence classification with
a \texttt{[CLS]} (or pooled) vector and a small linear head with four
logits. Sentiment on a review, spam versus ham on an email, and
``is this medical abstract about oncology?'' are the same shape: one
document in, one class out. The homework's AG~News fine-tune of TinyBERT
or MobileBERT is this application in miniature.

\paragraph{One label per word.}
Named-entity recognition---mark which tokens are people, places, or
organizations---is sequence labeling. So is part-of-speech tagging. The
encoder still runs once over the sentence; the head is applied at every
position. Intuition: the word ``Washington'' needs both left and right
context (``\ldots\ visited Washington last\ldots'' versus ``Washington
signed\ldots'') before the tag is reliable, which is why a bidirectional
encoder helps.

\paragraph{A span inside the input, not a new sentence.}
Extractive question answering---``when was so-and-so born?'' with the
answer already sitting in the paragraph as ``1953''---is span extraction.
The model does not invent the year; it points at the start and end tokens.
Highlighting the relevant sentence in a contract, or finding the price in
a product page, is the same idea. Generation would be needed only if the
answer had to be paraphrased or lengthened beyond a contiguous span.

\paragraph{Fine-tuning versus a frozen encoder.}
Two practical recipes use the same encoder-only idea.
\begin{itemize}[leftmargin=*]
  \item \textbf{Fine-tune the encoder and the head} (Part~1 of the
        homework). Start from a pretrained TinyBERT or MobileBERT, attach
        a four-way classification head, and update almost all weights on
        AG~News. The body adapts to news topics.
  \item \textbf{Freeze the encoder; train only a linear head} (Part~2).
        A sentence transformer turns each article into a fixed vector;
        logistic regression learns \(W\mathbf{e}+\mathbf{b}\) for the four
        topics. The body never moves. This is cheaper and works when the
        frozen embeddings already separate the classes; a raw encoder that
        was never trained as a sentence embedder often needs the fine-tune
        path instead.
\end{itemize}

\paragraph{What encoder-only is not for.}
Chat replies, translation into a differently ordered language, and
summaries that must invent new wording are length-changing generation
problems. Those need a decoder (or an encoder--decoder), which Lecture~9
takes up. If the answer must be a string that is not a contiguous piece of
the input, an encoder-only classifier or span head is the wrong tool.

\subsection{What the encoder is trained to predict}

The block does not choose the loss. The lecture trains this same encoder in
two regimes, and only one of them asks a person for labels.

The three tasks above are supervised. Someone has decided the class of the
sentence, the tag of each word, or the start and end of the answer. The
encoder's \(T\) vectors are the features, and a classifier is trained
against those decisions.

The other regime is the lecture's pre-training, developed fully with BERT
in the next part of the lecture \cite{devlin2019bert}. Mask a word that is
already in the sentence and ask the encoder to name it. ``the coffee was
\texttt{[MASK]}'' has target \texttt{cold}, copied out of the original
sentence before the mask was applied. The output at the mask position is a
classifier over the vocabulary, the same kind of cross-entropy as the
part-of-speech classifier, with one class per word type. The target was not
written by an annotator. The lecture calls the source unlabeled text,
meaning nobody was paid to label it, and calls this first stage
pre-training. The supervised tasks above are the fine-tuning stage: start
from those parameters, attach the complaint classifier or the tag
classifier, and continue training on the smaller labeled set.

So an encoder-only transformer is supervised when the target is a human
decision about the sentence, and the pre-training target is a word the
sentence already contained.

%----------------------------------------------------------------------
\begin{thebibliography}{9}
\bibitem{ba2016layer}
J.~L. Ba, J.~R. Kiros, and G.~E. Hinton.
Layer Normalization.
\emph{arXiv:1607.06450}, 2016.

\bibitem{ioffe2015batch}
S.~Ioffe and C.~Szegedy.
Batch Normalization: Accelerating Deep Network Training by Reducing
Internal Covariate Shift.
\emph{ICML}, 2015.

\bibitem{he2016deep}
K.~He, X.~Zhang, S.~Ren, and J.~Sun.
Deep Residual Learning for Image Recognition.
\emph{CVPR}, 2016.

\bibitem{vaswani2017attention}
A.~Vaswani, N.~Shazeer, N.~Parmar, J.~Uszkoreit, L.~Jones, A.~N. Gomez,
\L.~Kaiser, and I.~Polosukhin.
Attention Is All You Need.
\emph{NeurIPS}, 2017.

\bibitem{dao2022flashattention}
T.~Dao, D.~Y. Fu, S.~Ermon, A.~Rudra, and C.~R\'{e}.
FlashAttention: Fast and Memory-Efficient Exact Attention with IO-Awareness.
\emph{NeurIPS}, 2022.

\bibitem{shazeer2019fast}
N.~Shazeer.
Fast Transformer Decoding: One Write-Head is All You Need.
\emph{arXiv:1911.02150}, 2019.

\bibitem{ainslie2023gqa}
J.~Ainslie, J.~Lee-Thorp, M.~de Jong, Y.~Zemlyanskiy, F.~Lebr\'{o}n,
and S.~Sanghai.
GQA: Training Generalized Multi-Query Transformer Models from Multi-Head
Checkpoints.
\emph{EMNLP}, 2023.

\bibitem{devlin2019bert}
J.~Devlin, M.-W. Chang, K.~Lee, and K.~Toutanova.
BERT: Pre-training of Deep Bidirectional Transformers for Language
Understanding.
\emph{NAACL}, 2019.
\end{thebibliography}

\end{document}
